\section{Architecture et déploiement en production}

L'environnement de développement sécurisé présenté en section 3.4 pose les fondations de sécurité nécessaires. Cette section détaille l'architecture de production qui expose le système aux utilisateurs finaux via une API REST, une interface web, et un pipeline de déploiement automatisé garantissant la disponibilité et la fiabilité du service.

\subsection{API REST de prédiction}

L'API de prédiction constitue le point d'entrée principal du système en production. Implémentée avec FastAPI et servie derrière un reverse proxy Nginx, elle expose deux modèles (baseline et sécurisé) avec des garanties de performance et de sécurité.

\paragraph{Architecture de l'API FastAPI}

L'application principale (\texttt{src/api/app.py}) implémente une API asynchrone avec FastAPI et Uvicorn comme serveur ASGI. Au démarrage, l'API charge les deux modèles PyTorch et initialise les middlewares de sécurité :

\begin{lstlisting}[language=Python, caption=Initialisation de l'API FastAPI]
app = FastAPI(
    title="Secure Object Detection API",
    version="1.0.0",
    description="API de detection d'objets dangereux vs surs"
)

# Chargement des modeles au demarrage
baseline_model = torch.load('models/baseline/best_model.pth')
secured_model = load_encrypted_model('models/secured/best_secured_model.pth')

# Device detection : CUDA si disponible, sinon CPU
device = torch.device('cuda' if torch.cuda.is_available() else 'cpu')

# Middlewares de securite (Zone 4)
app.add_middleware(SecurityMiddleware)
app.add_middleware(PrometheusInstrumentator)
\end{lstlisting}

Le chargement du modèle sécurisé tente d'abord un déchiffrement AES-256-GCM (Zone 1), puis bascule sur le fichier non chiffré en cas d'absence de clé. La détection automatique du device (CUDA/CPU) optimise les performances selon le hardware disponible, avec un fallback CPU garantissant le fonctionnement sur toute infrastructure.

\paragraph{Endpoint de prédiction}

L'endpoint principal \texttt{POST /predict/\{model\_type\}} accepte une image en multipart form et retourne la prédiction avec métriques de confiance :

\begin{lstlisting}[language=Python, caption=Endpoint de prédiction]
@app.post("/predict/{model_type}")
async def predict(
    model_type: str,
    file: UploadFile = File(...),
    current_user: dict = Depends(verify_token)
):
    # Validation WAF (Zone 4)
    waf_check = waf.check_request(request.client.host, file.filename)
    if not waf_check['allowed']:
        raise HTTPException(status_code=429, detail=waf_check['reason'])

    # Preprocessing : resize + normalize ImageNet
    image = Image.open(file.file).convert('RGB')
    image_tensor = preprocess(image).unsqueeze(0).to(device)

    # Inference
    model = baseline_model if model_type == 'baseline' else secured_model
    with torch.no_grad():
        output = model(image_tensor)
        confidence, prediction = torch.max(F.softmax(output, dim=1), 1)

    # Audit logging (Zone 5)
    audit_id = audit_logger.log_event(
        event_type='PREDICTION',
        details={'model': model_type, 'prediction': class_names[prediction]},
        user_id=current_user['username'],
        image=image_tensor.cpu().numpy()
    )

    return {
        'model': model_type,
        'prediction': class_names[prediction.item()],
        'confidence': confidence.item(),
        'audit_id': audit_id
    }
\end{lstlisting}

Le pipeline intègre systématiquement les trois zones de sécurité : validation WAF (Zone 4) pour bloquer les requêtes malveillantes, preprocessing standardisé (resize 224x224, normalisation ImageNet) pour garantir la cohérence des entrées, et audit immutable (Zone 5) traçant chaque prédiction avec hachage de l'image et pseudonymisation de l'utilisateur.

\paragraph{Reverse proxy Nginx et load balancing}

Nginx (\texttt{docker/nginx.conf}) agit comme reverse proxy devant les deux instances d'API (ports 8001 et 8002), routant les requêtes selon le paramètre \texttt{model\_type} :

\begin{lstlisting}[language=nginx, caption=Configuration Nginx (extrait)]
upstream baseline_backend {
    server baseline-api:8000;
}

upstream secured_backend {
    server secured-api:8000;
}

server {
    listen 80;
    server_name localhost;

    location /predict/baseline {
        proxy_pass http://baseline_backend;
        proxy_set_header X-Real-IP $remote_addr;
        proxy_set_header X-Forwarded-For $proxy_add_x_forwarded_for;
    }

    location /predict/secured {
        proxy_pass http://secured_backend;
        proxy_set_header X-Real-IP $remote_addr;
        proxy_set_header X-Forwarded-For $proxy_add_x_forwarded_for;
    }

    # CORS headers pour interface web
    add_header Access-Control-Allow-Origin *;
    add_header Access-Control-Allow-Methods 'GET, POST, OPTIONS';
}
\end{lstlisting}

La configuration préserve l'IP réelle du client via les headers \texttt{X-Real-IP} et \texttt{X-Forwarded-For}, essentiels pour le rate limiting WAF par IP. Les headers CORS autorisent les requêtes cross-origin depuis l'interface web hébergée sur un domaine différent. Nginx compresse également les réponses avec gzip pour optimiser la bande passante.

\paragraph{Métriques Prometheus}

L'API expose un endpoint \texttt{/metrics} au format Prometheus pour le monitoring temps réel (Zone 5) :

\begin{lstlisting}[language=Python, caption=Métriques Prometheus exportées]
# Compteurs de requetes
http_requests_total = Counter(
    'http_requests_total',
    'Total HTTP requests',
    ['method', 'endpoint', 'status']
)

# Histogramme de latence
processing_time_seconds = Histogram(
    'processing_time_seconds',
    'Request processing time',
    ['model_type']
)

# Gauge d'accuracy
ai_model_accuracy = Gauge(
    'ai_model_accuracy',
    'Model accuracy',
    ['model_type']
)

# Compteur de predictions
ai_model_predictions_total = Counter(
    'ai_model_predictions_total',
    'Total predictions',
    ['model_type', 'prediction']
)
\end{lstlisting}

Ces métriques permettent à Grafana de construire des dashboards temps réel affichant le débit de requêtes, la distribution des latences (P50/P95/P99), l'accuracy par modèle, et la répartition des prédictions safe/dangerous. Les alertes Prometheus (configurées dans \texttt{configs/monitoring/alert\_rules.yml}) se déclenchent sur des seuils critiques : taux d'erreur >10\%, latence P95 >2s, ou accuracy <80\%.


\subsection{Interface utilisateur web et dashboard}

L'interface web (\texttt{web/index.html}) fournit un portail utilisateur pour soumettre des images, visualiser les prédictions, et consulter l'audit trail. Un dashboard administratif (\texttt{web/admin.html}) expose les métriques de sécurité de la Zone 4.

\paragraph{Interface de prédiction}

L'interface principale propose un formulaire de soumission d'image avec sélection du modèle :

\begin{lstlisting}[language=HTML, caption=Interface de prédiction (extrait)]
<div class="upload-zone">
    <input type="file" id="imageInput" accept="image/*">
    <select id="modelSelect">
        <option value="baseline">Baseline (non securise)</option>
        <option value="secured">Secured (adversarial training)</option>
    </select>
    <button onclick="submitPrediction()">Analyser l'image</button>
</div>

<div id="results" class="hidden">
    <h3>Resultat : <span id="prediction"></span></h3>
    <p>Confiance : <span id="confidence"></span>%</p>
    <p>Modele : <span id="model"></span></p>
    <p>ID Audit : <span id="auditId"></span></p>
</div>
\end{lstlisting}

La fonction JavaScript \texttt{submitPrediction()} envoie une requête POST multipart à l'API Nginx (port 80), reçoit la réponse JSON, et affiche dynamiquement le résultat avec un code couleur (vert pour safe, rouge pour dangerous). L'ID d'audit permet de retrouver la prédiction dans les logs pour conformité RGPD.

\paragraph{Dashboard administratif}

Le dashboard admin intègre les métriques de sécurité provenant des endpoints \texttt{/security/*} :

\begin{lstlisting}[language=JavaScript, caption=Chargement des metriques de securite]
async function loadSecurityDashboard() {
    // WAF status
    const wafStatus = await fetch('/security/waf/status', {
        headers: {'Authorization': `Bearer ${token}`}
    }).then(r => r.json());

    document.getElementById('blockedIPs').textContent = wafStatus.blocked_ips.length;
    document.getElementById('totalRequests').textContent = wafStatus.total_requests;

    // Anomalies recentes
    const anomalies = await fetch('/security/anomalies/recent', {
        headers: {'Authorization': `Bearer ${token}`}
    }).then(r => r.json());

    displayAnomalies(anomalies);

    // Refresh toutes les 5 secondes
    setTimeout(loadSecurityDashboard, 5000);
}
\end{lstlisting}

Le dashboard affiche le nombre d'IPs bloquées par le WAF, les anomalies récentes détectées (Zone 4), le taux d'erreur API, et les événements de sécurité (attaques détectées, échecs d'authentification). Un refresh automatique toutes les 5 secondes garantit la visibilité temps réel des menaces.

\paragraph{Authentification JWT dans l'interface}

L'interface web stocke le token JWT dans \texttt{localStorage} après login et l'inclut dans tous les appels API :

\begin{lstlisting}[language=JavaScript, caption=Gestion de l'authentification]
// Login
async function login(username, password) {
    const response = await fetch('/security/login', {
        method: 'POST',
        headers: {'Content-Type': 'application/json'},
        body: JSON.stringify({username, password})
    });

    const data = await response.json();
    localStorage.setItem('token', data.access_token);
    localStorage.setItem('role', data.role);

    // Redirection selon role
    if (data.role === 'admin') {
        window.location.href = '/admin.html';
    } else {
        window.location.href = '/index.html';
    }
}

// Verification du token a chaque requete
function getAuthHeaders() {
    const token = localStorage.getItem('token');
    return {'Authorization': `Bearer ${token}`};
}
\end{lstlisting}

Le rôle stocké (\texttt{admin}, \texttt{agent}, \texttt{guest}) conditionne l'affichage des menus : seuls les admins voient le lien vers le dashboard de sécurité, les agents accèdent à l'audit trail complet, et les guests ne peuvent que soumettre des prédictions. Cette séparation front-end est renforcée côté API par la vérification RBAC (Zone 4).


\subsection{Configurations de déploiement}

Trois configurations Docker Compose sont proposées pour adapter le déploiement aux contraintes de ressources et aux besoins de monitoring.

\paragraph{Configuration MINIMAL}

Le fichier \texttt{docker-compose.prod.MINIMAL.yml} déploie uniquement les composants essentiels sans base de données ni cache :

\begin{lstlisting}[language=yaml, caption=Configuration MINIMAL (extrait)]
services:
  baseline-api:
    build:
      context: .
      dockerfile: docker/Dockerfile.baseline
    ports:
      - "8001:8000"
    environment:
      - MODEL_TYPE=baseline
    networks:
      - secure-ai-network

  secured-api:
    build:
      context: .
      dockerfile: docker/Dockerfile.secured
    ports:
      - "8002:8000"
    environment:
      - MODEL_TYPE=secured
      - ENABLE_DEFENSES=true

  nginx:
    image: nginx:alpine
    ports:
      - "80:80"
    volumes:
      - ./docker/nginx.conf:/etc/nginx/nginx.conf:ro

  prometheus:
    image: prom/prometheus:latest
    ports:
      - "9090:9090"
    volumes:
      - ./configs/monitoring/prometheus.yml:/etc/prometheus/prometheus.yml

  grafana:
    image: grafana/grafana:latest
    ports:
      - "3000:3000"
\end{lstlisting}

Cette configuration (5 services) convient aux environnements de test ou aux déploiements à faible charge. Sans PostgreSQL, l'audit trail s'écrit uniquement en JSONL sur disque. Sans Redis, le WAF stocke les compteurs de rate limiting en mémoire (perdus au redémarrage). La latence d'inférence reste identique, seules les fonctionnalités d'audit persistant sont réduites.

\paragraph{Configuration STANDARD}

Le fichier \texttt{docker-compose.prod.yml} ajoute PostgreSQL et Redis pour un déploiement production standard :

\begin{lstlisting}[language=yaml, caption=Configuration STANDARD - Ressources (extrait)]
services:
  secured-api:
    deploy:
      resources:
        limits:
          cpus: '2.0'
          memory: 4G
        reservations:
          cpus: '1.0'
          memory: 2G
    restart: unless-stopped

  postgres:
    image: postgres:15-alpine
    environment:
      POSTGRES_USER: ${POSTGRES_USER}
      POSTGRES_PASSWORD: ${POSTGRES_PASSWORD}
      POSTGRES_DB: ai_metrics
    volumes:
      - postgres-data-prod:/var/lib/postgresql/data
    ports:
      - "127.0.0.1:5432:5432"
    healthcheck:
      test: ["CMD-SHELL", "pg_isready -U ${POSTGRES_USER}"]
      interval: 10s
      timeout: 5s
      retries: 5

  redis:
    image: redis:7-alpine
    command: redis-server --requirepass ${REDIS_PASSWORD} --appendonly yes
    volumes:
      - redis-data-prod:/data
    ports:
      - "127.0.0.1:6379:6379"

volumes:
  postgres-data-prod:
  redis-data-prod:
\end{lstlisting}

Les limites de ressources (2 CPU, 4GB pour le modèle sécurisé) préviennent l'épuisement mémoire lors de pics de charge. La politique \texttt{restart: unless-stopped} garantit le redémarrage automatique après une panne système. PostgreSQL et Redis exposent leurs ports uniquement sur localhost (\texttt{127.0.0.1}), isolant la base du réseau externe. Redis utilise AOF (Append-Only File) pour la persistance, reconstituant l'état du cache après redémarrage.

\paragraph{Configuration COMPLETE}

Le fichier \texttt{docker-compose.prod.COMPLETE.yml} inclut Loki/Promtail pour l'agrégation de logs et AlertManager pour les notifications :

\begin{lstlisting}[language=yaml, caption=Configuration COMPLETE - Monitoring avance]
  loki:
    image: grafana/loki:latest
    ports:
      - "3100:3100"
    volumes:
      - ./docker/loki/loki-config.yml:/etc/loki/loki-config.yml
      - loki-data:/loki

  promtail:
    image: grafana/promtail:latest
    volumes:
      - ./logs:/logs
      - ./docker/promtail/promtail-config.yml:/etc/promtail/config.yml
    command: -config.file=/etc/promtail/config.yml

  alertmanager:
    image: prom/alertmanager:latest
    ports:
      - "9093:9093"
    volumes:
      - ./docker/alertmanager/alertmanager.yml:/etc/alertmanager/alertmanager.yml
    command:
      - '--config.file=/etc/alertmanager/alertmanager.yml'
\end{lstlisting}

Loki centralise les logs applicatifs et système dans un format queryable depuis Grafana, alternative légère à Elasticsearch. Promtail scrape le répertoire \texttt{logs/} et transmet les entrées à Loki. AlertManager reçoit les alertes Prometheus et les route vers des webhooks externes (Slack, email, PagerDuty). Cette configuration (11 services) convient aux environnements critiques nécessitant observabilité complète et alerting 24/7.


\subsection{Pipeline CI/CD et tests automatisés}

Le pipeline de déploiement continu (\texttt{.github/workflows/deploy-production.yml}) automatise les tests de sécurité, la construction des images, et le déploiement avec rollback automatique en cas d'échec.

\paragraph{Workflow GitHub Actions}

Le workflow se déclenche automatiquement sur chaque push vers \texttt{main} et exécute 6 jobs séquentiels :

\begin{lstlisting}[language=yaml, caption=Pipeline CI/CD (extrait)]
name: Deploy to Production

on:
  push:
    branches: [main]
    paths-ignore: ['**.md', 'docs/**']
  workflow_dispatch:
    inputs:
      environment:
        description: 'Environment to deploy'
        required: true
        default: 'production'

jobs:
  security-tests:
    runs-on: ubuntu-latest
    steps:
      - name: Run unit tests
        run: docker-compose -f docker/docker-compose.dev.yml run --rm api pytest

      - name: Security scan with Bandit
        run: |
          pip install bandit
          bandit -r src/ -ll

      - name: Dependency audit
        run: |
          pip install safety
          safety check --json

  build-production:
    needs: security-tests
    steps:
      - name: Build baseline image
        run: docker build -f docker/Dockerfile.baseline -t baseline-api:latest .

      - name: Build secured image
        run: docker build -f docker/Dockerfile.secured -t secured-api:latest .

      - name: Test critical imports
        run: |
          docker run baseline-api:latest python -c "import fastapi, torch"
          docker run secured-api:latest python -c "import opacus"

  deploy:
    needs: build-production
    steps:
      - name: Deploy with docker-compose
        run: docker-compose -f docker/docker-compose.prod.yml up -d

      - name: Wait for health checks
        run: |
          timeout 60 bash -c 'until curl -f http://localhost:8001/health; do sleep 2; done'
          timeout 60 bash -c 'until curl -f http://localhost:8002/health; do sleep 2; done'
\end{lstlisting}

Le job \texttt{security-tests} exécute d'abord les tests unitaires dans un container isolé, puis scanne le code avec Bandit (détection de vulnérabilités Python niveau LL : high severity uniquement), et audite les dépendances avec Safety pour identifier les CVE connues. Le job \texttt{build-production} construit les deux images Docker et vérifie que les imports critiques (fastapi, torch, opacus) fonctionnent. Le job \texttt{deploy} lance le stack complet et attend jusqu'à 60 secondes que les endpoints \texttt{/health} répondent.

\paragraph{Tests post-déploiement}

Le job \texttt{smoke-tests} valide le déploiement avec des requêtes réelles :

\begin{lstlisting}[language=yaml, caption=Smoke tests]
  smoke-tests:
    needs: deploy
    steps:
      - name: Test baseline API
        run: |
          curl -f http://localhost:8001/health || exit 1
          curl -X POST http://localhost:8001/predict/baseline \
            -F "file=@tests/fixtures/sample_safe.jpg" | grep -q "safe"

      - name: Test secured API
        run: |
          curl -f http://localhost:8002/health || exit 1
          curl -X POST http://localhost:8002/predict/secured \
            -F "file=@tests/fixtures/sample_dangerous.jpg" | grep -q "dangerous"

      - name: Test web interface
        run: curl -f http://localhost:8080 || exit 1

      - name: Test Prometheus metrics
        run: curl -f http://localhost:9090/-/healthy || exit 1

      - name: Test Grafana
        run: curl -f http://localhost:3000/api/health || exit 1
\end{lstlisting}

Ces tests vérifient que chaque composant répond correctement : les deux APIs acceptent des images et retournent des prédictions cohérentes (safe pour \texttt{sample\_safe.jpg}, dangerous pour \texttt{sample\_dangerous.jpg}), l'interface web est accessible, et le stack de monitoring (Prometheus, Grafana) fonctionne. Un échec sur l'un de ces tests déclenche le job de rollback.

\paragraph{Rollback automatique}

Le job \texttt{rollback} s'exécute uniquement en cas d'échec des tests (\texttt{if: failure()}) :

\begin{lstlisting}[language=yaml, caption=Rollback automatique]
  rollback:
    needs: [deploy, smoke-tests]
    if: failure()
    steps:
      - name: Stop failed deployment
        run: docker-compose -f docker/docker-compose.prod.yml down

      - name: Clean containers
        run: docker system prune -f

      - name: Notify failure
        run: |
          echo "Deployment failed, rolled back"
          # Integration Slack/email ici
\end{lstlisting}

Le rollback arrête immédiatement tous les containers, nettoie les ressources Docker (images inutilisées, volumes anonymes), et envoie une notification d'échec. Ce mécanisme garantit qu'un déploiement défaillant n'impacte jamais la production : soit les tests passent et le nouveau code est déployé, soit ils échouent et l'état précédent est restauré.

\paragraph{Makefile pour déploiements manuels}

En complément du CI/CD, le Makefile fournit des commandes pour déploiements manuels et maintenance :

\begin{lstlisting}[language=bash, caption=Commandes Makefile (extrait)]
# Deploiement production complet avec rebuild
make prod

# Deploiement rapide sans rebuild
make prod-fast

# Tests post-deploiement
make test-prod

# Health checks
make health-prod

# Logs en temps reel
make logs-prod

# Statut des containers
make status-prod

# Arret propre
make stop-prod
\end{lstlisting}

La commande \texttt{make prod} reconstruit toutes les images (utile après modification du code), tandis que \texttt{make prod-fast} redémarre les containers avec les images existantes (déploiement en <30 secondes). La commande \texttt{make health-prod} vérifie que tous les services répondent, facilitant le diagnostic rapide après déploiement.